\section{Experimental Procedure}
\label{sec:exp_procedure}

This section details the experimental protocol that implements the CREATE framework described in Section~\ref{sec:method}. We describe the setup, the six conditions under comparison, and the training and evaluation pipeline. Results are deferred to Section~\ref{sec:results}.

\subsection{Experimental Setup and Configuration}
\label{sec:exp_setup}

Experiments are conducted on two Gymnasium environments. LunarLander-v3 serves as the primary benchmark for baseline comparison and mechanism ablation; BipedalWalker-v3 provides a secondary validation on continuous-action locomotion. Every reward program trains a freshly initialized PPO policy through Stable-Baselines3. LunarLander uses $10^6$ environment steps per candidate and BipedalWalker uses $5\times10^6$, with four vectorized environments inside each policy-training job. The computational infrastructure comprises two NVIDIA GeForce RTX 3090 GPUs (24\,GB each), with DeepSeek V4 Pro serving as the LLM backbone for all LunarLander conditions. Full prompts, model settings, hyperparameters, and per-seed traces are provided in the supplementary material.

\subsection{Condition Definitions}
\label{sec:conditions}

Six conditions are compared on LunarLander-v3 under a shared LLM backbone, environment description, reward interface, code checks, and PPO budget. Each condition contains five independent reward-search lineages and at most ten complete reward evaluations, so that every condition operates under the same total policy-training cost.

\textbf{Single-shot generation} evaluates only the initial reward produced by CREATE's task-understanding phase (Section~\ref{sec:init}), without subsequent observation or revision.

\textbf{Independent-10} generates ten unrelated initial reward candidates, each from an independent invocation of the reward-initialization phase, and evaluates each with a separate policy-training run. The best candidate is selected post hoc. This condition spends the same maximum number of evaluations as CREATE but without cross-round information flow.

\textbf{Stateless iterative revision} repeatedly revises the current reward from basic current-round feedback (native fitness, episode length, termination statistics, component episode sums), but removes the evidence organizer, intervention memory, cross-round component comparison, and hierarchical semantic editing.

\textbf{CREATE w/o evidence enrichment} retains the full loop, memory, best archive, and hierarchical editing, but removes the tool-using evidence organizer: the editor receives native fitness, episode length, termination statistics, and component episode sums, but lacks activation rates, magnitude shares, temporal trends, and cross-round summaries.

\textbf{CREATE w/o hierarchical editing} retains enriched evidence, memory, and the best archive, but removes the semantic-localization constraint (Eq.~\ref{eq:semantic_action}), allowing unconstrained multi-component rewrites.

\textbf{CREATE} runs the complete architecture described in Section~\ref{sec:method}.

\subsection{Training and Evaluation Protocol}
\label{sec:training_protocol}

Each reward-design round consists of a complete PPO training run followed by native evaluation. The policy is optimized with the generated reward $R_t$, whereas candidate quality is measured only by the untouched environment reward $r_{\mathrm{env}}$ (Eq.~\ref{eq:train_eval}). We define \emph{search fitness} as the mean undiscounted native episodic return over 20 fixed evaluation episodes and \emph{best search fitness} as the largest value retained within a lineage. After reward selection, the returned policy is evaluated over 100 unseen episodes; this is reported as \emph{test fitness}. The task-solving thresholds are 200 for LunarLander and 300 for BipedalWalker.

Between rounds, the evidence organizer (Eq.~\ref{eq:agent_loop}, $\Phi_{\mathcal T}$) compresses training records into a structured summary. The reward-editing agent then selects one principal semantic, chooses an edit level (L1/L2/L3), and applies a bounded modification (Eq.~\ref{eq:semantic_action}). The modified program is validated before the next round. The best archive is updated only when a new candidate exceeds the current incumbent in native search fitness (Eq.~\ref{eq:best_reward}). On BipedalWalker-v3, only the complete CREATE workflow is evaluated; no baseline or ablation comparison is conducted on that task.
